%% Generated by Sphinx.
\def\sphinxdocclass{report}
\documentclass[letterpaper,10pt,english]{sphinxmanual}
\ifdefined\pdfpxdimen
   \let\sphinxpxdimen\pdfpxdimen\else\newdimen\sphinxpxdimen
\fi \sphinxpxdimen=.75bp\relax
\ifdefined\pdfimageresolution
    \pdfimageresolution= \numexpr \dimexpr1in\relax/\sphinxpxdimen\relax
\fi
\newdimen\sphinxremdimen\sphinxremdimen = 10pt
%% let collapsible pdf bookmarks panel have high depth per default
\PassOptionsToPackage{bookmarksdepth=5}{hyperref}

\PassOptionsToPackage{booktabs}{sphinx}
\PassOptionsToPackage{colorrows}{sphinx}

\PassOptionsToPackage{warn}{textcomp}
\usepackage[utf8]{inputenc}
\ifdefined\DeclareUnicodeCharacter
% support both utf8 and utf8x syntaxes
  \ifdefined\DeclareUnicodeCharacterAsOptional
    \def\sphinxDUC#1{\DeclareUnicodeCharacter{"#1}}
  \else
    \let\sphinxDUC\DeclareUnicodeCharacter
  \fi
  \sphinxDUC{00A0}{\nobreakspace}
  \sphinxDUC{2500}{\sphinxunichar{2500}}
  \sphinxDUC{2502}{\sphinxunichar{2502}}
  \sphinxDUC{2514}{\sphinxunichar{2514}}
  \sphinxDUC{251C}{\sphinxunichar{251C}}
  \sphinxDUC{2572}{\textbackslash}
\fi
\usepackage{cmap}
\usepackage[T1]{fontenc}
\usepackage{amsmath,amssymb,amstext}
\usepackage{babel}


\usepackage{mathptmx}\usepackage[scaled=.9]{helvet}\usepackage{courier}



\usepackage{sphinx}

\fvset{fontsize=auto}
\usepackage{geometry}


% Include hyperref last.
\usepackage{hyperref}
% Fix anchor placement for figures with captions.
\usepackage{hypcap}% it must be loaded after hyperref.
% Set up styles of URL: it should be placed after hyperref.
\urlstyle{same}

\addto\captionsenglish{\renewcommand{\contentsname}{Tutorials}}

\usepackage{sphinxmessages}


\newenvironment{bcwchunk}[2]{\def\FrameCommand{{\color{#1}\vrule width 3pt}\fboxsep=6pt\colorbox{#2}}\MakeFramed{\advance\hsize-\width\FrameRestore}}{\endMakeFramed}
\definecolor{bcwbargoal}{HTML}{7B3FA0}
\definecolor{bcwbackgoal}{HTML}{F7F7F7}
\newenvironment{sphinxclassgoal}{\begin{bcwchunk}{bcwbargoal}{bcwbackgoal}}{\end{bcwchunk}}
\definecolor{bcwbarrequirement}{HTML}{1F5FBF}
\definecolor{bcwbackrequirement}{HTML}{EDF3FC}
\newenvironment{sphinxclassrequirement}{\begin{bcwchunk}{bcwbarrequirement}{bcwbackrequirement}}{\end{bcwchunk}}
\definecolor{bcwbarparameter}{HTML}{0E7C86}
\definecolor{bcwbackparameter}{HTML}{F7F7F7}
\newenvironment{sphinxclassparameter}{\begin{bcwchunk}{bcwbarparameter}{bcwbackparameter}}{\end{bcwchunk}}
\definecolor{bcwbardefinition}{HTML}{2E7D32}
\definecolor{bcwbackdefinition}{HTML}{EDF6EE}
\newenvironment{sphinxclassdefinition}{\begin{bcwchunk}{bcwbardefinition}{bcwbackdefinition}}{\end{bcwchunk}}
\definecolor{bcwbarrationale}{HTML}{8A8A8A}
\definecolor{bcwbackrationale}{HTML}{F7F7F7}
\newenvironment{sphinxclassrationale}{\begin{bcwchunk}{bcwbarrationale}{bcwbackrationale}}{\end{bcwchunk}}
\definecolor{bcwbardiscussion}{HTML}{8A8A8A}
\definecolor{bcwbackdiscussion}{HTML}{F7F7F7}
\newenvironment{sphinxclassdiscussion}{\begin{bcwchunk}{bcwbardiscussion}{bcwbackdiscussion}}{\end{bcwchunk}}
\definecolor{bcwbartarget}{HTML}{C26A00}
\definecolor{bcwbacktarget}{HTML}{F7F7F7}
\newenvironment{sphinxclasstarget}{\begin{bcwchunk}{bcwbartarget}{bcwbacktarget}}{\end{bcwchunk}}
\definecolor{bcwbaropen}{HTML}{C26A00}
\definecolor{bcwbackopen}{HTML}{F7F7F7}
\newenvironment{sphinxclassopen}{\begin{bcwchunk}{bcwbaropen}{bcwbackopen}}{\end{bcwchunk}}


\title{Book}
\date{}
\release{}
\author{Author}
\newcommand{\sphinxlogo}{\vbox{}}
\renewcommand{\releasename}{}
\makeindex
\begin{document}

\ifdefined\shorthandoff
  \ifnum\catcode`\=\string=\active\shorthandoff{=}\fi
  \ifnum\catcode`\"=\active\shorthandoff{"}\fi
\fi

\pagestyle{empty}
\sphinxmaketitle
\pagestyle{plain}
\renewcommand{\contentsname}{Contents}\sphinxtableofcontents
\pagestyle{normal}
\phantomsection\label{\detokenize{index::doc}}


\part*{Tutorials}\addcontentsline{toc}{part}{Tutorials}

\sphinxstepscope


\chapter{Guide}
\label{\detokenize{guide/guide:guide}}\label{\detokenize{guide/guide::doc}}

\section{Overview}
\label{\detokenize{guide/guide:overview}}
\sphinxAtStartPar
Text.

\part*{Reference}\addcontentsline{toc}{part}{Reference}

\sphinxstepscope


\chapter{Design}
\label{\detokenize{design/design:design}}\label{\detokenize{design/design::doc}}

\section{Overview}
\label{\detokenize{design/design:overview}}
\sphinxAtStartPar
Text.


\section{Goals}
\label{\detokenize{design/design:goals}}
\begin{sphinxuseclass}{chunk}
\begin{sphinxuseclass}{goal}\phantomsection\label{\detokenize{design/design:design.timing}}
\sphinxAtStartPar
\sphinxstylestrong{GOAL} \sphinxcode{\sphinxupquote{design.timing}}

\sphinxAtStartPar
No thread can change the timing.

\end{sphinxuseclass}
\end{sphinxuseclass}
\sphinxstepscope


\chapter{Core}
\label{\detokenize{core/core:core}}\label{\detokenize{core/core::doc}}

\section{Overview}
\label{\detokenize{core/core:overview}}
\sphinxAtStartPar
The core runs every thread through one pipeline.


\section{Rotation}
\label{\detokenize{core/core:rotation}}
\begin{sphinxuseclass}{chunk}
\begin{sphinxuseclass}{requirement}\phantomsection\label{\detokenize{core/core:core.rotation}}
\sphinxAtStartPar
\sphinxstylestrong{REQUIREMENT} \sphinxcode{\sphinxupquote{core.rotation}}

\sphinxAtStartPar
The core shall give the turn after
thread \sphinxstyleemphasis{t} to thread \sphinxstyleemphasis{t} + 1.

\sphinxAtStartPar
\sphinxstyleemphasis{Formal twin}

\begingroup\fvset{fontsize=\small}

\begin{sphinxVerbatim}[commandchars=\\\{\}]
\PYG{k}{def}\PYG{+w}{ }\PYG{n+nf}{core\PYGZus{}rotate}\PYG{p}{(}\PYG{n}{turn}\PYG{p}{)}\PYG{p}{:}
    \PYG{k}{return} \PYG{p}{\PYGZob{}}\PYG{l+s+s1}{\PYGZsq{}}\PYG{l+s+s1}{next}\PYG{l+s+s1}{\PYGZsq{}}\PYG{p}{:} \PYG{p}{(}\PYG{n}{turn} \PYG{o}{+} \PYG{l+m+mi}{1}\PYG{p}{)} \PYG{o}{\PYGZpc{}} \PYG{l+m+mi}{8}\PYG{p}{\PYGZcb{}}
\end{sphinxVerbatim}

\endgroup

\sphinxAtStartPar
Serves: {\hyperref[\detokenize{design/design:design.timing}]{\sphinxcrossref{\sphinxcode{\sphinxupquote{design.timing}}}}}

\end{sphinxuseclass}
\end{sphinxuseclass}
\sphinxAtStartPar
Verilog: {\hyperref[\detokenize{core/core:source-1}]{\sphinxcrossref{build/rtl/core/core\_rotate.v}}}

\sphinxAtStartPar
\sphinxstyleemphasis{Check (equiv)}: the module \sphinxcode{\sphinxupquote{core\_rotate}} equals the twin \sphinxcode{\sphinxupquote{core\_rotate}}.

\sphinxAtStartPar
Verifies: {\hyperref[\detokenize{core/core:core.rotation}]{\sphinxcrossref{\sphinxcode{\sphinxupquote{core.rotation}}}}}

\sphinxAtStartPar
Why: {\hyperref[\detokenize{core/core:why-1}]{\sphinxcrossref{see the explanation}}}

\begin{sphinxuseclass}{chunk}
\begin{sphinxuseclass}{open}
\sphinxAtStartPar
\sphinxstylestrong{OPEN} The thread count is not settled.

\end{sphinxuseclass}
\end{sphinxuseclass}
\begin{sphinxuseclass}{chunk}
\begin{sphinxuseclass}{definition}\phantomsection\label{\detokenize{core/core:core.turn}}
\sphinxAtStartPar
\sphinxstylestrong{DEFINITION} \sphinxcode{\sphinxupquote{core.turn}}

\sphinxAtStartPar
A \sphinxstyleemphasis{turn} is a thread’s cycle in the rotation.

\sphinxAtStartPar
Serves: {\hyperref[\detokenize{core/core:core.core}]{\sphinxcrossref{\sphinxcode{\sphinxupquote{core.core}}}}}

\end{sphinxuseclass}
\end{sphinxuseclass}
\begin{sphinxuseclass}{chunk}
\begin{sphinxuseclass}{definition}\phantomsection\label{\detokenize{core/core:core.core}}
\sphinxAtStartPar
\sphinxstylestrong{DEFINITION} \sphinxcode{\sphinxupquote{core.core}}

\sphinxAtStartPar
The \sphinxstyleemphasis{core} runs the threads in turn.

\sphinxAtStartPar
Serves: {\hyperref[\detokenize{design/design:design.timing}]{\sphinxcrossref{\sphinxcode{\sphinxupquote{design.timing}}}}}

\end{sphinxuseclass}
\end{sphinxuseclass}
\sphinxAtStartPar
\sphinxstylestrong{Thread.} An unlabelled bold paragraph is prose.


\section{Goals}
\label{\detokenize{core/core:goals}}
\begin{sphinxuseclass}{chunk}
\begin{sphinxuseclass}{goal}\phantomsection\label{\detokenize{core/core:core.timing}}
\sphinxAtStartPar
\sphinxstylestrong{GOAL} \sphinxcode{\sphinxupquote{core.timing}}

\sphinxAtStartPar
No thread can change the timing of another thread.

\end{sphinxuseclass}
\end{sphinxuseclass}
\begin{sphinxuseclass}{chunk}
\begin{sphinxuseclass}{parameter}\phantomsection\label{\detokenize{core/core:core.threads}}
\sphinxAtStartPar
\sphinxstylestrong{PARAMETER} \sphinxcode{\sphinxupquote{core.threads}} = 8 threads

\sphinxAtStartPar
The number of threads.

\sphinxAtStartPar
Serves: {\hyperref[\detokenize{core/core:core.core}]{\sphinxcrossref{\sphinxcode{\sphinxupquote{core.core}}}}}

\end{sphinxuseclass}
\end{sphinxuseclass}
\begin{sphinxuseclass}{chunk}
\begin{sphinxuseclass}{parameter}\phantomsection\label{\detokenize{core/core:core.turn-width}}
\sphinxAtStartPar
\sphinxstylestrong{PARAMETER} \sphinxcode{\sphinxupquote{core.turn\sphinxhyphen{}width}} = clog2(core.threads) = 3 bits

\sphinxAtStartPar
The width.

\sphinxAtStartPar
Serves: {\hyperref[\detokenize{core/core:core.threads}]{\sphinxcrossref{\sphinxcode{\sphinxupquote{core.threads}}}}}

\end{sphinxuseclass}
\end{sphinxuseclass}
\begin{sphinxuseclass}{chunk}
\begin{sphinxuseclass}{parameter}\phantomsection\label{\detokenize{core/core:core.spare}}
\sphinxAtStartPar
\sphinxstylestrong{PARAMETER} \sphinxcode{\sphinxupquote{core.spare}} = 2

\sphinxAtStartPar
A spare value.

\sphinxAtStartPar
Serves: {\hyperref[\detokenize{core/core:core.core}]{\sphinxcrossref{\sphinxcode{\sphinxupquote{core.core}}}}}

\end{sphinxuseclass}
\end{sphinxuseclass}
\begin{sphinxuseclass}{chunk}
\begin{sphinxuseclass}{target}\phantomsection\label{\detokenize{core/core:core.aim}}
\sphinxAtStartPar
\sphinxstylestrong{TARGET} \sphinxcode{\sphinxupquote{core.aim}} = core.threads * 2 = 16 threads

\sphinxAtStartPar
The aim of the core.

\sphinxAtStartPar
Serves: {\hyperref[\detokenize{core/core:core.threads}]{\sphinxcrossref{\sphinxcode{\sphinxupquote{core.threads}}}}}

\end{sphinxuseclass}
\end{sphinxuseclass}
\sphinxAtStartPar
Verilog: {\hyperref[\detokenize{core/core:source-2}]{\sphinxcrossref{build/rtl/core/pair.v}}}

\sphinxAtStartPar
Uses: {\hyperref[\detokenize{core/core:fragment-core.pair-logic}]{\sphinxcrossref{\sphinxcode{\sphinxupquote{:core.pair\sphinxhyphen{}logic}}}}}

\sphinxAtStartPar
Fragment: {\hyperref[\detokenize{core/core:source-3}]{\sphinxcrossref{:core.pair\sphinxhyphen{}logic}}}

\sphinxAtStartPar
\sphinxstyleemphasis{Check (test)}

\begingroup\fvset{fontsize=\small}

\begin{sphinxVerbatim}[commandchars=\\\{\}]
\PYG{k}{module}\PYG{+w}{ }\PYG{n}{tb}\PYG{p}{;}
\PYG{+w}{    }\PYG{k}{initial}\PYG{+w}{ }\PYG{n+nb}{\PYGZdl{}finish}\PYG{p}{;}
\PYG{k}{endmodule}
\end{sphinxVerbatim}

\endgroup

\sphinxAtStartPar
Verifies: {\hyperref[\detokenize{core/core:core.rotation}]{\sphinxcrossref{\sphinxcode{\sphinxupquote{core.rotation}}}}}

\sphinxAtStartPar
\sphinxstyleemphasis{Check (equiv)}: the module \sphinxcode{\sphinxupquote{core\_rotate}} equals the twin \sphinxcode{\sphinxupquote{core\_next}}.

\sphinxAtStartPar
Verifies: {\hyperref[\detokenize{core/core:core.rotation}]{\sphinxcrossref{\sphinxcode{\sphinxupquote{core.rotation}}}}}


\section{Explanation}
\label{\detokenize{core/core:explanation-1}}
\begin{sphinxuseclass}{explanation}\phantomsection\label{\detokenize{core/core:why-1}}\subsubsection*{{\hyperref[\detokenize{core/core:rotation}]{\sphinxcrossref{Rotation}}}}

\begin{sphinxuseclass}{chunk}
\begin{sphinxuseclass}{rationale}
\sphinxAtStartPar
\sphinxstylestrong{RATIONALE}

\sphinxAtStartPar
The aim is {\hyperref[\detokenize{core/core:core.aim}]{\sphinxcrossref{\DUrole{param}{16 threads}}}}. A thread’s instructions are eight cycles apart, one for each of {\hyperref[\detokenize{core/core:core.threads}]{\sphinxcrossref{\DUrole{param}{8 threads}}}}.

\end{sphinxuseclass}
\end{sphinxuseclass}
\end{sphinxuseclass}

\section{Implementation}
\label{\detokenize{core/core:implementation-1}}
\begin{sphinxuseclass}{implementation}\phantomsection\label{\detokenize{core/core:source-1}}\subsubsection*{Verilog: build/rtl/core/core\_rotate.v}

\def\sphinxLiteralBlockLabel{\label{\detokenize{core/core:file-build/rtl/core/core_rotate.v}}}
\begin{sphinxVerbatim}[commandchars=\\\{\}]
\PYG{k}{module}\PYG{+w}{ }\PYG{n}{core\PYGZus{}rotate}\PYG{+w}{ }\PYG{p}{(}\PYG{k}{input}\PYG{+w}{ }\PYG{k+kt}{wire}\PYG{+w}{ }\PYG{p}{[}\PYG{l+m+mh}{2}\PYG{o}{:}\PYG{l+m+mh}{0}\PYG{p}{]}\PYG{+w}{ }\PYG{n}{turn}\PYG{p}{,}\PYG{+w}{ }\PYG{k}{output}\PYG{+w}{ }\PYG{k+kt}{wire}\PYG{+w}{ }\PYG{p}{[}\PYG{l+m+mh}{2}\PYG{o}{:}\PYG{l+m+mh}{0}\PYG{p}{]}\PYG{+w}{ }\PYG{n}{next}\PYG{p}{)}\PYG{p}{;}
\PYG{+w}{    }\PYG{k}{assign}\PYG{+w}{ }\PYG{n}{next}\PYG{+w}{ }\PYG{o}{=}\PYG{+w}{ }\PYG{n}{turn}\PYG{+w}{ }\PYG{o}{+}\PYG{+w}{ }\PYG{l+m+mh}{3}\PYG{l+m+mi}{\PYGZsq{}d1}\PYG{p}{;}
\PYG{k}{endmodule}
\end{sphinxVerbatim}

\end{sphinxuseclass}
\begin{sphinxuseclass}{implementation}\phantomsection\label{\detokenize{core/core:source-2}}\subsubsection*{Verilog: build/rtl/core/pair.v}

\def\sphinxLiteralBlockLabel{\label{\detokenize{core/core:file-build/rtl/core/pair.v}}}
\begin{sphinxVerbatim}[commandchars=\\\{\}]
\PYG{k}{module}\PYG{+w}{ }\PYG{n}{pair}\PYG{+w}{ }\PYG{p}{(}\PYG{k}{input}\PYG{+w}{ }\PYG{k+kt}{wire}\PYG{+w}{ }\PYG{n}{a}\PYG{p}{,}\PYG{+w}{ }\PYG{k}{output}\PYG{+w}{ }\PYG{k+kt}{wire}\PYG{+w}{ }\PYG{n}{b}\PYG{p}{)}\PYG{p}{;}
\PYG{+w}{    }\PYG{o}{\PYGZlt{}}\PYG{o}{\PYGZlt{}}\PYG{o}{:}\PYG{n}{core}\PYG{p}{.}\PYG{n}{pair}\PYG{o}{\PYGZhy{}}\PYG{k+kt}{logic}\PYG{o}{\PYGZgt{}}\PYG{o}{\PYGZgt{}}
\PYG{k}{endmodule}
\end{sphinxVerbatim}

\end{sphinxuseclass}
\begin{sphinxuseclass}{implementation}\phantomsection\label{\detokenize{core/core:source-3}}\subsubsection*{Fragment: :core.pair\sphinxhyphen{}logic}

\def\sphinxLiteralBlockLabel{\label{\detokenize{core/core:fragment-core.pair-logic}}}
\begin{sphinxVerbatim}[commandchars=\\\{\}]
assign b = a;
\end{sphinxVerbatim}

\end{sphinxuseclass}\subsubsection*{Index of code}
\begin{itemize}
\item {} 
\sphinxAtStartPar
\sphinxcode{\sphinxupquote{build/rtl/core/core\_rotate.v}}: defined in {\hyperref[\detokenize{core/core:file-build/rtl/core/core_rotate.v}]{\sphinxcrossref{Core}}}.

\item {} 
\sphinxAtStartPar
\sphinxcode{\sphinxupquote{build/rtl/core/pair.v}}: defined in {\hyperref[\detokenize{core/core:file-build/rtl/core/pair.v}]{\sphinxcrossref{Core}}}.

\item {} 
\sphinxAtStartPar
\sphinxcode{\sphinxupquote{:core.pair\sphinxhyphen{}logic}}: defined in {\hyperref[\detokenize{core/core:fragment-core.pair-logic}]{\sphinxcrossref{Core}}}, used in {\hyperref[\detokenize{core/core:file-build/rtl/core/pair.v}]{\sphinxcrossref{\sphinxcode{\sphinxupquote{build/rtl/core/pair.v}}}}}.

\end{itemize}



\renewcommand{\indexname}{Index}
\printindex
\end{document}